%======================================================================
\section{Conclusões}\label{sec:conclusoes}
%======================================================================
\subsection{Situação}

O sistema está projetado de ponta a ponta e, na maior parte, implementado
(Tabela~\ref{tab:objetivos}). Os sensores foram escolhidos e caracterizados; o condicionamento
do LM35 foi projetado, simulado e validado, e na placa nova passou a funcionar com uma
alimentação só e um filtro na entrada do ADC; a PCB única foi desenhada, verificada e
fabricada; o firmware das duas placas está completo e compila sem avisos; a detecção de queda
passou de um limiar a um modelo embarcado; e a central funciona com um simulador no lugar do
hub. O que falta é sobretudo bancada: montar a placa e verificar, com o hardware real, o que
hoje está verificado só no computador.

\begin{table}[H]
\centering
\caption{Objetivos específicos e situação.}
\label{tab:objetivos}
\small
\begin{tabular}{@{}>{\raggedright\arraybackslash}p{6.4cm}>{\raggedright\arraybackslash}p{7.4cm}@{}}
\toprule
Objetivo & Situação \\
\midrule
Selecionar e caracterizar os sensores & feito (primeiro semestre); a faixa do MPU passou a $\pm$8\,g, confirmada na placa \\
Condicionar o LM35 & feito e validado em protoboard; falta conferir na PCB v1 \\
PCB única & desenhada, verificada (ERC e DRC) e fabricada; falta montar e testar \\
Rede CAN & firmware pronto e validado com o professor numa versão anterior; falta o teste entre as placas v1 \\
Interface sem fio e contrato MQTT & firmware pronto, contrato documentado; falta o teste ESP-NOW em bancada \\
Detecção de queda embarcada & modelo no firmware, com paridade C $\times$ Python; provisório até os dados do pulso \\
Regras e alertas no hub & prontas e testadas no computador \\
Histórico, painel e notificação & funcionando com o simulador; falta o bot do Telegram com o celular do cuidador \\
\bottomrule
\end{tabular}
\end{table}

\subsection{Limitações}

\begin{itemize}
  \item O modelo de queda foi treinado com o sensor na cintura, e o vestível vai no pulso.
  \item Nada do sistema completo foi verificado em bancada ainda: as metas de latência,
        disponibilidade e autonomia estão definidas, mas não medidas.
  \item A exatidão da temperatura na placa nova ainda não foi conferida contra um termopar.
  \item O vestível usa a mesma placa do hub: é grande e depende de um \textit{power bank}, sem
        medição de bateria.
  \item A cobertura do ESP-NOW fica limitada ao raio do hub; o corredor e as áreas comuns
        estão fora do escopo desta fase.
\end{itemize}

\subsection{Próximos passos}

\begin{enumerate}
  \item Montar as placas v1 e fazer os testes da Seção~\ref{sec:verificacao}, começando pelo
        ESP-NOW entre o vestível e o hub, para a apresentação do Q3 (24/11).
  \item Coletar quedas e atividades com o vestível no pulso, com várias pessoas, e retreinar
        o modelo; reavaliar ali os atributos de frequência e o SVM.
  \item Medir o consumo e implementar o \textit{light sleep} no vestível.
  \item Uma placa própria para o vestível, com bateria LiPo, carregador e medição da bateria.
  \item Na central: um servidor NTP local, para a ILPI funcionar sem internet; um painel para
        a telemetria do vestível; e um banco relacional (PostgreSQL) com residentes,
        dispositivos, cuidadores e o ciclo de vida dos alertas.
\end{enumerate}
